% from ega4-9-fr.tex
\oldpage[IV-3]{72}

supposer que $f$ est lui-même un morphisme propre~; on sait alors qu'il en est
de même de $f_s$ pour tout $s\in S$ (\textbf{II}, 5.4.2, (iii)). Supposons au
contraire que $\eta\in S-E$, et distinguons deux cas~:

\begin{enumerate}[label=\upshape{$\arabic{enumi}^\circ$}]
  \item Supposons que $f_\eta$ soit non séparé~; alors il résulte de (iii) que
  $f_s$ est non séparé (et \emph{a fortiori} non propre) dans un voisinage de
  $\eta$.

  \item Supposons $f_\eta$ séparé~; $Y$ est réunion finie d'ouverts affines
  $V_i$, et pour que $f_s$ soit propre, il faut et il suffit que chacune de ses
  restrictions $f_s^{-1}((V_i)_s)\to(V_i)_s$ le soit (\textbf{II}, 5.4.1)~;
  on peut donc se borner au cas où $Y$ est \emph{affine}, donc un schéma. Dire
  que $f_\eta$ n'est pas propre signifie (\textbf{II}, 5.6.3) qu'il existe un
  morphisme \emph{de type fini} $h:Z\to Y_\eta$ tel que le morphisme
  $(f_\eta)_{(Z)}:X_\eta\times_{Y_\eta}Z\to Z$ ne soit pas fermé. Comme $Y$ est
  un schéma, on déduit de (8.8.2, (i) et (ii)) (en restreignant au besoin $S$ à
  un voisinage de $\eta$) qu'il existe un morphisme de type fini $g:Y'\to Y$
  tel que $Z$ soit isomorphe à $Y'_\eta$, et que $g_\eta=h$~; si l'on pose
  $X'=X\times_YY'$, $f'=f_{(Y')}:X'\to Y'$, on a
  $(f_\eta)_{(Z)}=f'_\eta$, et par hypothèse il existe donc une partie fermée
  $M'$ de $X'_\eta$ telle que $f'_\eta(M')$ ne soit pas fermé dans $Y'_\eta$.
  Or $M'$ est la trace sur $X'_\eta$ d'une partie fermée $N'$ de $X'$~; comme
  $X'$ est noethérien et $f'$ de type fini, $f'(N')$ est constructible dans
  $Y'$ (1.8.4) et par hypothèse
  $(f'(N'))_\eta=f'_\eta(N'_\eta)=f'_\eta(M')$ n'est pas fermé dans
  $Y'_\eta$. On conclut alors de (9.5.4) qu'il existe un voisinage $U$ de
  $\eta$ dans $S$ tel que pour $s\in U$,
  $(f'(N'))_s=f'_s(N'_s)$ ne soit pas fermé dans $Y'_s$~; autrement dit, le
  morphisme $f'_s$ n'est pas fermé, et \emph{a fortiori} le morphisme $f_s$
  n'est pas propre.
\end{enumerate}

(vi) La propriété est conjonction des propriétés (iv) et (vii) (8.11.1), donc
la proposition résulte dans ce cas de ce qui a déjà été prouvé.

(ix) La vérification de la condition (9.2.1, (i)) résulte de (2.7.1, (xi)). On
peut donc se borner au cas où $S$ est affine, noethérien et intègre de point
générique $\eta$, et à prouver que $E$ ou $S-E$ est un voisinage de $\eta$. Si
$\eta\in E$, on peut (en remplaçant $S$ par un voisinage de $\eta$) supposer
que $f$ est une immersion fermée en vertu de (8.10.5, (iv)), et alors il est
clair que $f_s$ est une immersion fermée pour tout $s\in S$ (\textbf{I},
4.3.2). Supposons donc que $\eta\in S-E$ et distinguons deux cas~:

\begin{enumerate}[label=\upshape{$\arabic{enumi}^\circ$}]
  \item $f_\eta$ n'est pas un morphisme fini. Alors il résulte de (vi) que dans
  un voisinage de $\eta$, $f_s$ n'est pas fini, ni \emph{a fortiori} une
  immersion fermée.

  \item $f_\eta$ est fini~; alors, en vertu de (8.10.5, (x)), on peut supposer
  (en restreignant au besoin $S$) que $f$ lui-même est un morphisme \emph{fini}.
  Dans ce cas $\mathcal A=\mathcal A(X)=f_*(\mathcal O_X)$ est un
  $\mathcal O_Y$-Module cohérent (\textbf{II}, 6.1.3) et
  $f=\mathcal A(u)$, où $u:\mathcal O_Y\to\mathcal A$ est un homomorphisme de
  $\mathcal O_Y$-Algèbres (\textbf{II}, 1.1.2)~; comme $f_\eta$ n'est pas une
  immersion fermée par hypothèse, $u_\eta$ n'est pas surjectif (\textbf{II},
  1.4.10), donc (9.4.5) il existe un voisinage de $\eta$ dans lequel $u_s$
  n'est pas surjectif, et par suite (\textbf{I}, 4.2.3) $f_s$ n'est pas une
  immersion fermée.
\end{enumerate}

(viii) La vérification de la condition (9.2.1, (i)) se fait comme dans (ix), en
utilisant cette fois (2.7.1, (x)) et le fait que toute immersion d'un préschéma
noethérien dans un autre est quasi-compacte. On est donc ramené au cas où $S$
est affine, noethérien et intègre de point générique $\eta$, et à prouver que
$E$ ou $S-E$ est un voisinage de $\eta$. Si $\eta\in E$, on conclut comme dans
(ix) à l'aide de (8.10.5, (ii)). Si $\eta\in S-E$, on distingue de nouveau
deux cas~:

\begin{enumerate}[label=\upshape{$\arabic{enumi}^\circ$}]
  \item $f_\eta(X_\eta)$ n'est pas une partie localement fermée de $Y_\eta$.
  Comme $f(X)$ est constructible dans $Y$ (1.8.4) et que
  $f_s(X_s)=(f(X))_s$, on déduit de (9.5.4) que pour $s$ voisin de $\eta$,
  $f_s(X_s)$ n'est pas localement fermé dans $Y_s$, et \emph{a fortiori}
  $f_s$ n'est pas une immersion.

  \item